\documentclass[UTF8,fontset=fandol,11pt]{ctexart}
\usepackage[a4paper,margin=25mm]{geometry}
\usepackage{amsmath,amssymb,amsthm,mathtools}
\usepackage{xcolor}
\usepackage{booktabs,array,enumitem,fancyhdr,hyperref}
\hypersetup{colorlinks=true,linkcolor=blue!50!black,urlcolor=blue!50!black}
\setlength{\emergencystretch}{2em}
\setlength{\parskip}{0.3em}
\linespread{1.16}
\allowdisplaybreaks
\newtheorem{theorem}{定理}[section]
\newtheorem{proposition}[theorem]{命题}
\newtheorem{lemma}[theorem]{引理}
\theoremstyle{definition}
\newtheorem{definition}[theorem]{定义}
\theoremstyle{remark}
\newtheorem{remark}[theorem]{说明}
\newcommand{\dd}{\,\mathrm d}
\newcommand{\ii}{\mathrm i}
\newcommand{\Rea}{\operatorname{Re}}
\newcommand{\Ima}{\operatorname{Im}}
\newcommand{\Li}{\operatorname{Li}}
\newcommand{\supp}{\operatorname{supp}}
\pagestyle{fancy}
\fancyhf{}
\fancyhead[L]{\small 学习笔记}
\fancyhead[R]{\small 初稿}
\fancyfoot[C]{\thepage}
\setlength{\headheight}{15pt}
\author{}
\date{2026年10月22日初稿}
\makeatletter
\renewcommand{\maketitle}{%
  \thispagestyle{plain}%
  \begin{center}
    {\LARGE\bfseries \@title\par}
    \vspace{0.6em}
    {\small \@date\par}
  \end{center}
  \vspace{0.6em}}
\makeatother
\newcommand{\draftnotice}{\noindent{\small
本文为供逐篇审阅的学习笔记初稿，不是新增研究成果。
文中区分直接证明与引用结果；正文不依赖本文。
本轮已作公式、条件及排版初步核查，尚待作者逐句审定。}\par\medskip}

\usepackage{tikz}
\usetikzlibrary{arrows.meta}
\date{2026年10月22日初稿；10月23日增补示意图}
\title{Mellin 变换与对数坐标}
\begin{document}
\maketitle
\draftnotice
\section{为什么出现两种看起来不同的变换}
Fourier 变换以加法变量的振荡为基础；Mellin 变换中则出现 $x^{s-1}$。
二者的关系不是只靠形式相似：取对数变量，可以把正实轴上的乘法结构变为实轴上的加法结构。
本篇完整写出变量、测度与归一化的变化，并说明何时可以反演。
所需背景为实积分、复指数与 Fourier 反演定理。

图~\ref{fig:log-coordinates} 先列出 $u=\log x$ 的坐标关系。
当 $x$ 乘以固定常数 $a>0$ 时，$u$ 增加 $\log a$；同时
$\dd x/x=\dd u$，因而乘法卷积中的测度也变为加法变量的通常测度。
\begin{figure}[htbp]
\centering
\begin{tikzpicture}[x=1cm,y=1cm,>=Stealth,font=\small]
\draw[->] (0,2)--(13,2) node[right] {$x$};
\draw[->] (0,0)--(13,0) node[right] {$u$};
\foreach \posX/\posU/\tickUp/\tickDown in {
  1.125/1.5/{\frac14}/{-2\log2},
  1.75/4/{\frac12}/{-\log2},
  3/6.5/{1}/{0},
  5.5/9/{2}/{\log2},
  10.5/11.5/{4}/{2\log2}} {
  \draw[gray!65,dashed] (\posX,1.92)--(\posU,0.08);
  \draw (\posX,1.92)--(\posX,2.08) node[above=3pt] {$\tickUp$};
  \draw (\posU,0.08)--(\posU,-0.08) node[below=3pt] {$\tickDown$};
}
\node[fill=white,inner sep=3pt] at (8.5,1.05) {$u=\log x$};
\draw[->,blue!65!black,thick] (3,2.95)--(5.5,2.95)
  node[midway,above=2pt] {$x\mapsto2x$};
\draw[->,blue!65!black,thick] (6.5,-0.95)--(9,-0.95)
  node[midway,below=2pt] {$u\mapsto u+\log2$};
\end{tikzpicture}

\caption{正实轴与对数坐标的对应。上、下两行各自采用线性刻度；
虚线连接同一取值在两种坐标中的位置。$x$ 的等比取点在 $u$ 轴上等距，
伸缩 $x\mapsto2x$ 对应平移 $u\mapsto u+\log2$。}
\label{fig:log-coordinates}
\end{figure}

\section{固定一条竖线后作变量代换}
定义
\[
 \mathcal Mf(s)=\int_0^\infty f(x)x^{s-1}\dd x.
\]
固定实数 $c$，假设 $\int_0^\infty |f(x)|x^{c-1}\dd x<\infty$。
令 $x=e^u$，则 $\dd x=e^u\dd u$，并定义 $g_c(u)=e^{cu}f(e^u)$。
于是
\[
 \int_{\mathbb R}|g_c(u)|\dd u
 =\int_0^\infty|f(x)|x^{c-1}\dd x<\infty.
\]
采用 Fourier 惯例
$\widehat g(t)=\int_{\mathbb R}g(u)e^{-\ii tu}\dd u$，逐步计算得
\begin{align*}
 \mathcal Mf(c+\ii t)
 &=\int_{-\infty}^{\infty}f(e^u)e^{(c+\ii t-1)u}e^u\dd u\\
 &=\int_{\mathbb R}g_c(u)e^{\ii tu}\dd u
 =\widehat g_c(-t).
\end{align*}
这里有负号，是因为 Mellin 定义用了 $x^{+\ii t}$，而所选 Fourier 定义用了 $e^{-\ii tu}$。
改变 Fourier 惯例会改变符号和常数，不能混用。

加权因子 $e^{cu}$ 仍须保留，不能只按图~\ref{fig:log-coordinates} 换横坐标而略去它。

\section{从 Fourier 反演得到 Mellin 反演}
再假设 $\widehat g_c\in L^1(\mathbb R)$，且 $g_c$ 在所考察的 $u$ 处连续。
Fourier 反演给出
\[
 g_c(u)=\frac1{2\pi}\int_{\mathbb R}\widehat g_c(t)e^{\ii tu}\dd t
       =\frac1{2\pi}\int_{\mathbb R}\mathcal Mf(c+\ii t)e^{-\ii tu}\dd t.
\]
除以 $e^{cu}$ 并还原 $x=e^u$，得到
\[
 f(x)=\frac1{2\pi}\int_{\mathbb R}\mathcal Mf(c+\ii t)x^{-c-\ii t}\dd t
 =\frac1{2\pi\ii}\int_{c-\ii\infty}^{c+\ii\infty}
                 \mathcal Mf(s)x^{-s}\dd s.
\]
最后一步用了 $s=c+\ii t$、$\dd s=\ii\dd t$。
这组假设是方便使用的一组充分条件，不是所有反演形式的必要条件。
若变换不绝对可积，便须另行说明截断方式及反演定理，不能直接沿用本节证明。

\section{一个既能计算又能检验条件的例子}
令 $f(x)=(1-x)$（$0<x<1$），而 $x\ge1$ 时取 $0$。对 $\Rea s>0$，
\[
 \mathcal Mf(s)=\int_0^1(x^{s-1}-x^s)\dd x
 =\frac1s-\frac1{s+1}=\frac1{s(s+1)}.
\]
固定 $c>0$ 后，$g_c$ 绝对可积且连续；变换沿竖线为 $O((1+|t|)^{-2})$，也绝对可积。
因此上节反演适用，包括 $x=1$ 处，此时函数连续且值为零。
相反，$f(x)=1$（$0<x<1$）的变换为 $1/s$，沿竖线并不绝对可积。
阶跃核与连续的三角权重，正是在这里出现了反演条件的差别。

\section{乘法卷积为何化为乘积}
设 $f,g$ 在同一权重 $x^{c-1}$ 下绝对可积，定义
\[
 h(x)=\int_0^\infty f(x/y)g(y)\frac{\dd y}{y}
\]
（在积分存在的点取值）。先作绝对值估计，再令 $x=uy$：
\begin{align*}
 \int_0^\infty\int_0^\infty
 |f(x/y)g(y)|x^{c-1}\frac{\dd y}{y}\dd x
 &=\left(\int_0^\infty |f(u)|u^{c-1}\dd u\right)
   \left(\int_0^\infty |g(y)|y^{c-1}\dd y\right)<\infty.
\end{align*}
Tonelli 与 Fubini 定理因而保证几乎处处的存在性和换序合法性。
对 $s=c+\ii t$，同一代换给出
$\mathcal Mh(s)=\mathcal Mf(s)\mathcal Mg(s)$。
测度 $\dd y/y$ 使乘法伸缩保持自然形式；它并非可以随意换成 $\dd y$ 的装饰。

\section{与 Laplace 变换及振荡的联系}
若 $f$ 支持于 $(0,1]$，另取变量 $v=-\log x$，即 $x=e^{-v}$，则
\[
 \mathcal Mf(s)=\int_0^\infty f(e^{-v})e^{-sv}\dd v,
\]
在绝对收敛的区域内，这就是函数 $v\mapsto f(e^{-v})$ 的 Laplace 变换。
回到前面的 $u=\log x$，等式 $x^{\beta+\ii\gamma}=e^{\beta u}e^{\ii\gamma u}$ 表明：
幂函数在对数坐标中分解为包络与定频振荡。
这种代换解释了显式公式中的对数尺度频率，但并不将零点频率变成等间距 Fourier 级数。

\section*{参考依据与后续核验}
Fourier 反演作为标准定理引用；本文其余恒等式均在所列绝对可积条件下推导。
可参看本书《积分变换》及《零点项的对数尺度振荡》，不依赖它们的具体节号。
图~\ref{fig:log-coordinates} 的绘图源码随本文提供，坐标变换不依赖数值近似。
改变 $c$ 后加权函数形状的数值比较仍可另作补充，不影响本篇现有推导。
\end{document}
